% §4 Experimental Setup

\section{Experimental Setup}
\label{sec:setup}

We instantiate the measurement framework (\S\ref{sec:method}) on one locked $(\text{pool}, \text{dataset}, \text{protocol})$ configuration selected by requirements-first screening (\S\ref{sec:setup:config}).
All reported characterization and routing metrics use ARC-Challenge test ($n{=}1{,}172$); calibration ($n{=}299$) supports one-time $c(q)$ selection and nested model fitting only.
Studies follow a fixed hierarchy aligned with our prediction-then-decision pipeline: \textbf{characterization} (Studies I--II), \textbf{complementary predictive information} (Study III), and \textbf{decision utility} (Study IV).
Utility results are reported only after predictive information is established \citep{hu2024routerbench}.

\subsection{Configuration selection}
\label{sec:setup:config}

We select one primary configuration---a $(\text{pool}, \text{dataset}, \text{protocol})$ triple---using the following criteria:
\begin{enumerate}
  \item Objective correctness labels (exact letter match on multiple-choice items).
  \item A non-degenerate weak--strong outcome mix with sufficient routing opportunity.
  \item Compatible prefill probes: full-vocabulary logits at the probe step on every pool member.
  \item Feasible compute for offline oracle labeling and signal extraction.
\end{enumerate}

We screened candidate triples on small samples ($n \approx 10$--$50$, seed 42) to confirm prefill extraction succeeds and outcome buckets are non-degenerate.
Table~\ref{tab:setup} summarizes the locked configuration; the appendix reports excluded candidates (e.g., Qwen+GSM8K: 0\% opportunity).

\textbf{Locked configuration (2026-06-22):} Llama~3.2 1B / 3B Instruct + ARC-Challenge \citep{clark2018think}, prompt protocol v1, seed 42.

% T1: Locked experimental configuration
\begin{table}[t]
  \centering
  \small
  \begin{tabular}{ll}
    \toprule
    \textbf{Component} & \textbf{Locked value} \\
    \midrule
    Weak model ($M_w$)   & Llama-3.2-1B-Instruct \\
    Strong model ($M_s$) & Llama-3.2-3B-Instruct \\
    Primary dataset      & ARC-Challenge \citep{clark2018think} \\
    CALIB split          & Official validation ($n{=}299$) \\
    TEST split           & Official test ($n{=}1{,}172$) \\
    Prompt formatting     & Deterministic chat template (\texttt{prompt\_protocol.py}) \\
    Random seed          & 42 \\
    $c(q)$ selection     & \texttt{piece\_count} (CALIB screening) \\
    \bottomrule
  \end{tabular}
  \caption{Locked experimental configuration. Screening exclusions in appendix.}
  \label{tab:setup}
\end{table}


\subsection{Models and probe pool}
\label{sec:setup:models}

\begin{table}[t]
  \centering
  \small
  \begin{tabular}{ll}
    \toprule
    \textbf{Role} & \textbf{Model} \\
    \midrule
    Weak ($M_w$) & \texttt{meta-llama/Llama-3.2-1B-Instruct} \\
    Strong ($M_s$) & \texttt{meta-llama/Llama-3.2-3B-Instruct} \\
    \bottomrule
  \end{tabular}
  \caption{Locked model pool.}
  \label{tab:pool}
\end{table}

Both models share the Llama~3.2 tokenizer family.
We extract signals with Hugging Face Transformers; model revisions are pinned at experiment lock.
Full-vocabulary logprobs at prefill step $t=T$ are required for $H$ and $m$.

\subsection{Datasets and splits}
\label{sec:setup:splits}

\textbf{Primary: ARC-Challenge} \citep{clark2018think} --- multiple-choice science reasoning; gold letter match after task formatting (\S\ref{sec:method:protocol}).

\begin{table}[t]
  \centering
  \small
  \begin{tabular}{llrl}
    \toprule
    \textbf{Split role} & \textbf{Official HF split} & \textbf{$n$} & \textbf{Use} \\
    \midrule
    CALIB & \texttt{validation} & 299 & $c(q)$ calibration; router fitting \\
    TEST  & \texttt{test}       & 1{,}172 & All reported paper metrics \\
    \bottomrule
  \end{tabular}
  \caption{ARC-Challenge split policy. Train split is unused.}
  \label{tab:splits}
\end{table}

ARC train is unused.
We never tune on test oracle labels.

\textbf{Out-of-distribution generalization:} MMLU \citep{hendrycks2021measuring} --- three subjects, same protocol and frozen $c(q)$; one transfer table (\S\ref{sec:characterization}).
Calibration of $c(q)$ is not repeated on MMLU.

\subsection{Pre-study calibration of $c(q)$}
\label{sec:setup:calib}

Before main characterization, we select one representative $c(q)$ from five tokenizer-based candidates on calibration data only:

\begin{table}[t]
  \centering
  \small
  \begin{tabular}{ll}
    \toprule
    \textbf{Family} & \textbf{Candidate column} \\
    \midrule
    Length & \texttt{piece\_count} \\
    Lexical diversity & \texttt{mattr} \\
    Information & \texttt{text\_shannon}, \texttt{text\_shannon\_norm} \\
    Compressibility & \texttt{compression\_ratio} \\
    \bottomrule
  \end{tabular}
  \caption{Model-independent complexity candidates for $c(q)$ calibration.}
  \label{tab:cq-candidates}
\end{table}

For each candidate we compute Spearman $\rho$ and AUROC vs.\ $y^{\text{opp}}$ with bootstrap CIs.
Selection maximizes a composite of normalized $|\rho|$ and AUROC ranks; ties favor simpler features (length $\rightarrow$ compressibility $\rightarrow$ information $\rightarrow$ lexical diversity).
The winner is frozen as \texttt{c\_q} for all test analysis.
\textbf{[TBD: winning feature after screen run.]}

\subsection{Oracle labels}
\label{sec:setup:oracle}

For each $(q,M_i)$ on calibration and test splits:
\begin{enumerate}
  \item Build the same chat prompt $P(u)$ as signal extraction.
  \item Run greedy full inference (\texttt{max\_new\_tokens=8} for ARC letter answers).
  \item Parse output $\rightarrow$ $y(q,M_i)\in\{0,1\}$.
\end{enumerate}
Oracle labels define $y^{\text{opp}}$ and outcome buckets; they are never used to compute signals or to tune routing on the test set.

\begin{table}[t]
  \centering
  \small
  \begin{tabular}{lll}
    \toprule
    \textbf{Bucket} & \textbf{Condition} & \textbf{Interpretation} \\
    \midrule
    easy & $y_w{=}1$, $y_s{=}1$ & Weak model sufficient \\
    opportunity & $y_w{=}0$, $y_s{=}1$ & Escalation helps ($y^{\text{opp}}{=}1$) \\
    weak-only & $y_w{=}1$, $y_s{=}0$ & Strong model harmful \\
    too-hard & $y_w{=}0$, $y_s{=}0$ & Pool insufficient \\
    \bottomrule
  \end{tabular}
  \caption{Oracle outcome buckets.}
  \label{tab:buckets-def}
\end{table}

\subsection{Nested evaluation protocol}
\label{sec:setup:nested}

\begin{table}[t]
  \centering
  \small
  \begin{tabular}{lll}
    \toprule
    \textbf{Stage} & \textbf{Data} & \textbf{Uses oracle labels?} \\
    \midrule
    $c(q)$ calibration & CALIB & Yes (selection only) \\
    Study III logistic fit & CALIB & Yes (characterization) \\
    Study IV router fit + $\tau$ & CALIB & Yes (utility) \\
    Reported metrics & TEST & Evaluate only \\
    \bottomrule
  \end{tabular}
  \caption{Nested calibration/evaluation protocol.}
  \label{tab:nested}
\end{table}

Study III and Study IV fit separate logistic models on calibration data; Study IV does not reuse Study III weights.

\subsection{Metrics}
\label{sec:setup:metrics}

\begin{table}[t]
  \centering
  \small
  \begin{tabular}{llll}
    \toprule
    \textbf{Study} & \textbf{Layer} & \textbf{Question (RH)} & \textbf{Primary metrics} \\
    \midrule
    I  & Characterization & Does the model-independent family encode routing need? (RH1) &
      $\rho$, AUROC, AUPRC, bucket distributions \\
    II & Characterization & Does the model-dependent family encode routing need? (RH2) &
      Same (+ strong probes in appendix) \\
    III & Prediction & Complementary predictive information across families (RH3) &
      $\Delta$AUROC: $c \rightarrow c{+}H_w{+}m_w$; DeLong \\
    IV & Decision & Do predictions improve lightweight routing? (RH4) &
      Accuracy, avg.\ cost, utility, Pareto vs.\ baselines \\
    \bottomrule
  \end{tabular}
  \caption{Study hierarchy and metrics. All test metrics are evaluation-only.}
  \label{tab:metrics}
\end{table}

\textbf{Cost model:} weak $=1$, strong $=3$ (relative tier units).
Probe-cost analysis compares $m$ prefills vs.\ $m$ full generations \citep{guha2024smoothie,kotte2026ucci}.

\subsection{Study IV: utility and controllable trade-offs}
\label{sec:setup:utility}

Study IV evaluates \textbf{decision} utility after Study III establishes \textbf{prediction} quality.
We tune threshold $\tau$ on calibration data by maximizing
\begin{equation}
U = \text{accuracy} - \lambda \cdot \text{avg.\ cost},
\end{equation}
with $\lambda \in \{0,\,0.05,\,0.10\}$ representing deployment preferences (accuracy-first to cost-sensitive).
Each $\lambda$ yields one operating point; we report accuracy, average cost, utility, and a cost--quality Pareto plot (Figure~\ref{fig:cost-quality})---evaluation reporting borrowed from the routing literature \citep{hu2024routerbench,cao2026scope}, not a SCOPE-style learned estimator.
Study IV uses a separate CALIB-fit logistic model over $(c, H_w, m_w)$; it does not reuse Study III weights.

\subsection{Reproducibility}
\label{sec:setup:repro}

Fixed seed 42; split manifest \texttt{analysis/splits.json}; artifacts versioned under \texttt{experiments/}.
Model-independent features are extracted on CPU; oracle and probes on GPU with bfloat16 where supported.
